\bprob[title={DoCarmo, Riemannian Geometry, page 106 question 8}]

    Let $M$ be a connected Riemannian manifold of dimension $n\geq3$.
    Suppose that $M$ is {\emphcolor isotropic}: its sectional curvature at each point is independent on choice of subspace
    $\sigma\subseteq T_pM$.
    Prove that $M$ has constant sectional curvature; that is its sectional curvature is further independent on choice of point.

\eprob

Let us define the $4$-tensor
$$ \Rm'(W,Z,X,Y) = \iprod{W,X}\iprod{Z,Y} - \iprod{Z,X}\iprod{W,Y} . $$
Because sectional curvature $K(p,\sigma)=K(p)$ is independent of $\sigma$, we know that $\Rm=K\Rm'$.
Notice that for $U\in\X(M)$,
$$ \nabla_U\Rm'(W,Z,X,Y) = U(\Rm'(W,Z,X,Y)) - \iprod{\nabla_UW,X}\iprod{Z,Y} + \iprod{Z,X}\iprod{\nabla_UW,Y} - \cdots . $$
Because the connection is metric, further expansions of the terms show that this all cancels out.
This is lengthy so I will not write everything out, but for example expanding once gives
$$ U(\Rm'(W,Z,X,Y)) = U\iprod{W,X}\iprod{Z,Y} - U\iprod{Z,X}\iprod{W,Y} =
\iprod{\nabla_UW,Z}\iprod{Z,Y} - \iprod{Z,X}\iprod{\nabla_UW,Y} + \cdots $$
which cancels out with the first term above.

Thus
$$ \nabla_U\Rm' = 0 $$
and so
$$ \nabla\Rm(\bul,\bul,\bul,\bul,U) = \nabla_U\Rm = \nabla_U(K\Rm') = (UK)\Rm' + K\nabla_U\Rm' = (UK)\Rm' . $$
The second Binachi identity gives
$$ \nabla\Rm(W,Z,X,Y,U) + \nabla\Rm(W,Z,Y,U,X) + \nabla\Rm(W,Z,U,X,Y) = 0 . $$
Thus
$$ (UK)\Rm'(W,Z,X,Y) + (XK)\Rm'(W,Z,Y,U) + (YK)\Rm'(W,Z,U,X) = 0 . $$

Fix a $p\in M$ and let $X\in T_pM$ be nonzero.
Because $n\geq3$ we can find $Y,Z\in T_pM$ such that $\iprod{X,Y}=\iprod{Y,Z}=\iprod{Z,X}=0$ and $\iprod{Z,Z}=1$ (extend $X$
to an orthogonal basis).
Letting $U=Z$, and letting $W$ be arbitrary we get
$$ (ZK)\Rm'(W,Z,X,Y) + (XK)\Rm'(W,Z,Y,Z) + (YK)\Rm'(W,Z,Z,X) = 0 . $$
Expanding definitions,
$$ \eqalign{
    \Rm'(W,Z,X,Y) &= \iprod{W,X}\iprod{Z,Y} - \iprod{Z,X}\iprod{W,Y} = 0,\cr
    \Rm'(W,Z,Y,Z) &= \iprod{W,Y}\iprod{Z,Z} - \iprod{Z,Y}\iprod{W,Z} = \iprod{W,Y},\cr
    \Rm'(W,Z,Z,X) &= \iprod{W,Z}\iprod{Z,X} - \iprod{Z,Z}\iprod{W,X} = -\iprod{W,X}.\cr
} $$
Thus we see that
$$ (XK)\iprod{W,Y} - (YK)\iprod{W,X} = \iprod{W,(XK)Y - (YK)X} = 0 . $$
Since this is true for arbitrary $W$, this means that $(XK)Y-(YK)X=0$.
Because $X,Y$ are linearly independent this implies $XK=0$.

So for all $X\in T_pM$ we have $XK=0$ (for $X=0$ this is trivial).
Thus $K$ is locally constant, and since $M$ is connected this means $K$ is constant as desired.

\bprob[title={DoCarmo, Riemannian Geometry, page 187 question 10}]

    Let $f\colon\bar M\to M$ be a Riemannian submersion; $X,Y,Z,W\in\X(M)$ and $\bar X,\bar Y,\bar Z,\bar W$ be their
    horizontal lifts.
    Let $R,\bar R$ be the curvature tensors of $M,\bar M$ respectively.
    Prove the following.

    \benum
        \item
        $$ \iprod{\bar R(\bar X,\bar Y)\bar Z,\bar W} = \iprod{R(X,Y)Z,W} - \frac14\iprod{[\bar X,\bar Z]^v,[\bar Y,\bar W]^v}
        + \frac14\iprod{[\bar Y,\bar Z]^v,[\bar X,\bar W]^v} - \frac12\iprod{[\bar Z,\bar W]^v,[\bar X,\bar Y]^v} . $$
        \item Let $\sigma$ be the plane generated by orthonormal $X,Y\in\X(M)$, and $\bar\sigma$ the plane generated by
        $\bar X,\bar Y$.
        Then
        $$ K(\sigma) = \bar K(\bar\sigma) + \frac34\norm{[\bar X,\bar Y]^v}^2 \geq \bar K(\bar\sigma) . $$
    \eenum

\eprob

\benum
    \item Notice that, as per a previous exercise,
    $$ \bar X\iprod{\bar\nabla_{\bar Y}\bar Z,\bar W} = \bar X\iprod{\bar{\nabla_XY}+\frac12[\bar X,\bar Y]^v,\bar W}
    = \bar X\iprod{\bar{\nabla_XY},\bar W} = X\iprod{\nabla_XY,W} $$
    where the second equality is because $\bar W$ is horizontal and thus orthogonal to vertical vectors.

    Because the connection is metric,
    $$ \iprod{\bar\nabla_{\bar X}\bar\nabla_{\bar Y}\bar Z,\bar W} = \bar X\iprod{\bar\nabla_{\bar Y}\bar Z,\bar W}
    - \iprod{\bar\nabla_{\bar Y}\bar Z,\bar\nabla_{\bar X}\bar W} . $$
    The first term is equal to $X\iprod{\bar\nabla_YZ,W}$.
    The second term is equal to
    $$ \multline{
        \iprod{\bar\nabla_{\bar Y}\bar Z,\bar\nabla_{\bar X}\bar W} =
        \iprod{\bar{\nabla_YZ}+\frac12[\bar Y,\bar Z],\bar{\nabla_XW}+\frac12[\bar X,\bar W]^v}
        = \iprod{\bar{\nabla_YZ},\bar{\nabla_XW}} + \frac14\iprod{[\bar Y,\bar Z]^v,[\bar X,\bar W]^v}\cr
        = \iprod{\nabla_YZ,\nabla_XW} + \frac14\iprod{[\bar Y,\bar Z]^v,[\bar X,\bar W]^v} . 
    } $$
    Thus
    $$ \multline{
        \iprod{\bar\nabla_{\bar X}\bar\nabla_{\bar Y}\bar Z,\bar W} = X\iprod{\nabla_YZ,W} - \iprod{\nabla_YZ,\nabla_XW}
        - \frac14\iprod{[\bar Y,\bar Z]^v,[\bar X,\bar W]^v}\cr
        =\iprod{\nabla_X\nabla_YZ,W} - \frac14\iprod{[\bar Y,\bar Z]^v,[\bar X,\bar W]^v} .
    } $$

    If $T\in\X(\bar M)$ is vertical, then
    $$ \iprod{\bar\nabla_T\bar X,\bar Y} = \iprod{\bar\nabla_{\bar X}T,\bar Y} + \iprod{[T,\bar X],\bar Y} . $$
    The second term is zero, as per the previous exercise.
    Thus this is equal to (since $\iprod{T,\bar Y}=0$)
    $$ = \bar X\iprod{T,\bar Y} - \iprod{T,\bar\nabla_{\bar X}\bar Y} = -\iprod{T,\bar\nabla_{\bar X}\bar Y}
    = -\frac12\iprod{T,[\bar X,\bar Y]^v} . $$

    Now we have
    $$ \iprod{\bar\nabla_{[\bar X,\bar Y]}\bar Z,\bar W} = \iprod{\bar\nabla_{[\bar X,\bar Y]^h}\bar Z,\bar W}
    + \iprod{\bar\nabla_{[\bar X,\bar Y]^v}\bar Z,\bar W} . $$
    As we just showed, the second term is equal to
    $$ \iprod{\bar\nabla_{[\bar X,\bar Y]^v}\bar Z,\bar W} = -\frac12\iprod{[\bar X,\bar Y]^v,[\bar Z,\bar W]^v} . $$
    As per the previous homework $\iprod{[\bar X,\bar Y]^h,\bar Z}=\iprod{[\bar X,\bar Y],\bar Z}=\iprod{[X,Y],Z}$.
    Thus
    $$ \iprod{[\bar X,\bar Y]^h,\bar Z} = \iprod{\bar{[X,Y]},\bar Z} = \iprod{[X,Y],Z} $$
    and since all horizontal vectors are the lift of some vector field this means that $[\bar X,\bar Y]^h=\bar{[X,Y]}$.
    Thus
    $$ \iprod{\bar\nabla_{[\bar X,\bar Y]}\bar Z,\bar W} = \iprod{\nabla_{[X,Y]}Z,W} -
    \frac12\iprod{[\bar X,\bar Y]^v,[\bar Z,\bar W]^v} . $$

    Putting this together,
    $$ \kern\leftskip\eqalign{
        \iprod{\bar R(\bar X,\bar Y)\bar Z,\bar W} &=
        \iprod{\bar\nabla_{\bar Y}\bar\nabla_{\bar X}\bar Z,\bar W} -
        \iprod{\bar\nabla_{\bar X}\bar\nabla_{\bar Y}\bar Z,\bar W} + \iprod{\bar\nabla_{[\bar X,\bar Y]}\bar Z,\bar W}\cr
        &= \iprod{R(X,Y)Z,W} - \frac14\iprod{[\bar X,\bar Z]^v,[\bar Y,\bar W]^v}\cr
        &\hphantom{{}={}}+ \frac14\iprod{[\bar Y,\bar Z]^v,[\bar X,\bar W]^v} -
            \frac12\iprod{[\bar Z,\bar W]^v,[\bar X,\bar Y]^v} .
    } $$

    \item Because $df$ preserves the metric of horizontal vectors, $\bar X,\bar Y$ remain orthonormal.
    Thus
    $$ \bar K(\bar\sigma) = \bar\Rm(\bar X,\bar Y,\bar X,\bar Y) = \Rm(X,Y,X,Y) -
    \frac14\iprod{[\bar X,\bar X]^v,[\bar Y,\bar Y]^v} + \frac14\iprod{[\bar Y,\bar X]^v,[\bar X,\bar Y]^v}
    - \frac12\iprod{[\bar X,\bar Y]^v,[\bar X,\bar Y]^v} , $$
    and because the Lie bracket is antisymmetric we get that the first inner product is zero, the second is equal to
    $-\norm{[\bar X,\bar Y]^v}^2/4$ and so
    $$ K(\sigma) = \Rm(X,Y,X,Y) = \bar K(\bar\sigma) + \frac34\norm{[\bar X,\bar Y]^v}^2 . $$
\eenum

\bprob[title={DoCarmo, Riemannian Geometry, page 188 question 12}]

    Define a Riemannian metric on $\bC^{n+1}-0$ as follows.
    For $x\in\bC^{n+1}-0$ and $V,W\in T_z(\bC^{n+1}-0)\cong\bC^{n+1}$ define
    $$ \iprod{V,W}_x = \frac{\Re(V,W)}{(Z,Z)} . $$
    Notice that the metric $\iprod{\bul,\bul}$ restricted to $S^{2n+1}\subseteq\bC^{n+1}-0$ coincides with the one induced
    from $\bR^{2n+2}$.
    \benum
        \item Show that for all $0\leq\theta\leq2\pi$, $e^{i\theta}\colon S^{2n+1}\to S^{2n+1}$ is an isometry.
        Thus it is possible to define a Riemannian metric on $\CP^n$ such that $f\colon S^{2n+1}\to\CP^n$ is a Riemannian
        submersion.

        \item Show that in this metric, the sectional curvature of $\CP^n$ is
        $$ K(\sigma) = 1 + 3\cos^2\phi $$
        where $\sigma$ is generated by the orthonormal pair $X,Y$ and $\cos\phi=\iprod{\bar X,i\bar Y}$ where $\bar X$
        is the horizontal lift of $X$.
    \eenum

\eprob

\benum
    \item Multiplication by $z\in\bC$ has differential $dz\colon T_v\bC^n\to T_{zv}\bC^n$, after identifying both naturally with
    $\bC^n$, given by multiplication by $z$.
    Thus multiplication by a $\mu=e^{i\theta}$ satisfies for $V,W\in T_zS^{2n+1}$,
    $$ \iprod{d\mu(V),d\mu(W)} = \iprod{\mu V,\mu W} = \Re(\mu V,\mu W) = \Re\abs\mu^2(V,W) = \Re(V,W) = \iprod{V,W} . $$
    Thus $\mu$ is an isometry as desired.

    In order for $f$ to be a Riemannian submersion, the isomorphism
    $$ df_p\colon(T_pS^{2n+1})^h \xtos\sim T_{fp}\CP^n $$
    must preserve the metric.
    This uniquely defines the metric on $T_{fp}\CP^n$: define $\iprod{v,u}_{\CP^n}=\iprod{\bar v,\bar u}_{S^{2n+1}}$ where
    $\bar v,\bar u$ are the unique horizontal lifts of $v,u$.

    Now if $fp=fq$ then $q=\mu p$ for some $\mu=e^{i\theta}$ and because $f=f\circ\mu$,
    $$ df_q = df_p\circ d\mu_p $$
    so that if $df_q(v')=df_p(v)$ then as $\mu$ is an isometry, $v'=d\mu_p(v)$ and since $\mu$ is an isometry this means
    the metric on $\CP^n$ is well-defined.

    This requires that $\mu$ be an isometry between the horizontal subspaces, so that $d\mu_p$ maps horizontal vectors to
    horizontal vectors.
    Because horizontal vectors are those in the kernel of $df$, and since $f=f\circ\mu$, if $df(v)=0$ then
    $df\circ d\mu(v)=df(v)=0$ as desired.

    Note that the horizontal lift of a vector field $X\in\X(M)$ exists independently of whether or not $M$'s metric is smooth; as long as $f$ is a surjective submersion
    which maps horizontal vectors isometrically.
    And the horizontal lift is smooth (when $\bar M$'s metric is); the proof given in the previous homework doesn't have an assumption on $M$'s metric's smoothness.
    But then $g(X,Y)\circ f=\bar g(\bar X,\bar Y)$ is smooth (since $\bar g,\bar X,\bar Y$ are) and since $f$ is a surjective submersion, $g(X,Y)$ is smooth for all $X,Y\in\X(M)$
    so $g$ is.
    That is to say, if there is a surjective submersion $f\colon\bar M\to M$ between manifolds equipped with metrics, where $\bar M$'s is smooth, such that $f$ is an isometry on
    $\bar M$'s horizontal vectors (relative to $f$), then $M$'s metric is smooth.

    This is satisfied in our case, and so the metric defined on $\CP^n$ is smooth.
    Thus this defines a smooth metric on $\CP^n$ as desired.

    \item Let $Z$ be the position vector describing $S^{2n+1}$, i.e.\ $Z_v=v$ for $v\in S^{2n+1}$.
    Because its extension to $\bC^n$ is smooth, so $Z\in\Gamma(T\bC^n|_{S^{2n+1}})$.

    We know that $S^n$ is the level set of $1$ under $N\colon\bR^{n+1}-0\to\bR,N(x)=\norm x^2$ which is a submersion.
    Thus $T_pS^n$ consists of vectors $v\in T_p\bR^{n+1}$ where $dN(v)=0$.
    And the Jacobian of $N$ at $p$ is $2p$, so $dN(v)=2(p,v)$, i.e.\ $T_pS^n=p^\perp\subseteq\bR^{n+1}$.

    In particular, $ip\in T_pS^{2n+1}$ since
    $$ \iprod{p,ip}_p = \frac{\Re(p,ip)}{(p,p)} = \frac{-\Re i(p,p)}{(p,p)} = 0 . $$
    Thus $iZ$ (which maps $p\mapsto ip$) is tangent to $S^{2n+1}$, so it defines a smooth vector field $iZ\in\X(S^{2n+1})$.

    Let $X,Y\in\X(\CP^n)$, and let $\alpha\colon(-\epsilon,\epsilon)\to S^{2n+1}$ be a smooth curve with $\alpha(0)=p$ and $\alpha'(0)=\bar X(p)$.
    Let $\bar\nabla$ be the Levi-Civita connection of $S^{2n+1}$, then
    $$ (\bar\nabla_{\bar X}iZ)_p = \evdrv t_{t=0}iZ\circ\alpha(t) = \evdrv t_{t=0}i\alpha(t) = i\alpha'(0) = i\bar X(p) . $$
    That is to say,
    $$ \bar\nabla_{\bar X}iZ = i\bar X . $$

    Now,
    $$ \iprod{\bar\nabla_{\bar X}\bar Y,iZ} = \bar X\iprod{\bar Y,iZ} - \iprod{\bar Y,\bar\nabla_{\bar X}iZ} = -\iprod{\bar Y,i\bar X} , $$
    where the last equality is due to $\bar Y$ being horizontal and $iZ$ being vertical, as well as what we just showed.
    So we have
    $$ \iprod{[\bar X,\bar Y],iZ} = \iprod{\bar\nabla_{\bar X}\bar Y-\bar\nabla_{\bar Y}\bar X,iZ} = -\iprod{\bar Y,i\bar X} + \iprod{\bar X,i\bar Y} . $$
    By definition we have
    $$ \eqalign{
        \iprod{v,iu}_p &= \frac{\Re(v,iu)}{(p,p)} = \frac{-\Re i(v,u)}{(p,p)} = \frac{\Im (v,u)}{(p,p)},\cr
        \iprod{iv,u}_p &= -\frac{\Im(v,u)}{(p,p)},
    } $$
    so that $\iprod{iv,u}_p=-\iprod{v,iu}_p$.
    Thus
    $$ \iprod{[\bar X,\bar Y],iZ} = 2\iprod{\bar X,i\bar Y} = 2\cos\phi . $$

    By the previous question,
    $$ K(\sigma) = \bar K(\bar\sigma) + \frac34\norm{[\bar X,\bar Y]^v}^2 . $$
    Because the vertical tangent space has dimension $1$, it is spanned by $iZ$, which has unit norm.
    Thus
    $$ [\bar X,\bar Y]^v = \iprod{[\bar X,\bar Y],iZ}iZ = (2\cos\phi)iZ $$
    which has norm squared of $4\cos^2\phi$.
    Because the unit sphere has sectional curvature of $1$, we get that
    $$ K(\sigma) = 1 + 3\cos^2\phi $$
    as desired.
\eenum

\def\hess{{\rm Hess}}
\def\div{{\rm div}}

\bprob

    Let $U\subseteq\bR^n$ be open and $f\colon U\to\bR$ be smooth.
    Define
    $$ M = \set{(\bar x,y)}[\bar x\in U,\ y=f(\bar x)] \subseteq \bR^{n+1} $$
    and consider the coordinate chart $\x\colon U\to M$ by $\x(\bar x)=(\bar x,f(\bar x))$.
    Prove the following.
    \benum
        \item A unit normal field to $M$ is given by
        $$ \nu\circ\x = \frac{(-\nabla f,1)}{\sqrt{1+\norm{\nabla f}^2}} . $$
        \item The second fundamental form of $M$ is given by
        $$ \iprod{\sff(\partial_i,\partial_j),\nu}\circ\x = \frac{\hess f(\partial_i,\partial_j)}{\sqrt{1+\norm{\nabla f}^2}} . $$
        \item The mean curvature of $M$ is given by
        $$ \iprod{H,\nu}\circ\x = \div\parens{\frac{\nabla f}{\sqrt{1+\norm{\nabla f}^2}}} . $$
        \item If $U=\bR^n$ and $f$ is linear, then $M$ is a minimal surface.
        \item If $U=\set{u\in\bR^2}[\norm u>a]$ and $f=a\cosh^{-1}(\norm u/a)$ then $M$ is a minimal surface.
    \eenum

\eprob

\def\tpart{\tilde\partial}
\benum
    \item Define $F\colon U\times\bR\to\bR$ by $F(\bar x,y)=y-f(\bar x)$.
    Then $\nabla F=(-\nabla f,1)$ so that $F$ is a submersion and thus $M=F^{-1}0$ is embedded.
    And thus for $p\in M$, $T_pM=\ker\nabla F$ which are all the vectors orthogonal to $(-\nabla f,1)$.
    Normalizing this, we see that a normal vector field is
    $$ \frac{(-\nabla f,1)}{\sqrt{1+\norm{\nabla f}^2}} $$
    as desired.

    \item Let us denote by $\tilde\partial_i$ the coordinate frame induced by $(U,\x)$, and $\partial_i$ the Euclidean coordinate frame.
    We want to show
    $$ \iprod{\sff(\tilde\partial_i,\tilde\partial_j),\nu}\circ\x = \frac{\hess f(\partial_i,\partial_j)}{\sqrt{1+\norm{\nabla f}^2}} . $$
    Because $\sff(X,Y)=\bar\nabla_XY-\nabla_XY$ where $\bar\nabla$ is the Euclidean connection and $\nabla$ is $M$'s, since $\nabla_XY\in\X(M)$ is orthogonal to $\nu$,
    $$ \iprod{\sff(X,Y),\nu} = \iprod{\bar\nabla_XY,\nu} . $$
    From exercise 8 we know
    $$ \tilde\partial_i = f_i\partial_{n+1} + \partial_i,\qquad f_i=\pdvof f{\partial x_i} . $$
    And thus
    $$ \iprod{\sff(\tpart_i,\tpart_j),\nu} = \iprod{\bar\nabla_{\tpart_i}\tpart_j,\nu} = \iprod{\bar\nabla_{\partial_i+f_i\partial_{n+1}}(\partial_j+f_j\partial_{n+1}),\nu} . $$
    Because $\bar\nabla_{\partial_i}\partial_j=0$ for all $i,j$, we see that this is equal to
    $$ = \iprod{\bar\nabla_{\partial_i}(f_j\partial_{n+1})+f_i\bar\nabla_{\partial_{n+1}}(f_j\partial_{n+1}),\nu} = \iprod{f_{ij}\partial_{n+1},\nu} . $$
    Note the last equality is due to $f$ not being a function of $x_{n+1}$ so that $\partial_{n+1}f_j=0$.

    From the previous point, this is equal to
    $$ \iprod{\sff(\tpart_i,\tpart_j),\nu} = \frac{f_{ij}}{\sqrt{1+\norm{\nabla f}^2}} . $$

    Now we compute $\hess f(\partial_i,\partial_j)$ from definition:
    $$ \hess f(\partial_i,\partial_j) = \iprod{\bar\nabla_{\partial_i}\nabla f,\partial_j} = \iprod{\bar\nabla_{\partial_i}\sum_\ell f_\ell\partial_\ell,\partial_j}
    = \sum_\ell\iprod{f_{i\ell}\partial_\ell,\partial_j} = f_{ij} $$
    which gives us the desired result.

    \item \def\sqd{\sqrt{1+\norm{\nabla f}^2}} \def\den{1+\norm{\nabla f}^2}
    Firstly, mean curvature can be computed in the coordinate frame $\tpart_i$ by
    $$ H = g^{ij}\sff(\tpart_i,\tpart_j) , $$
    and so by the previous point
    $$ \iprod{H,\nu} = g^{ij}\iprod{\sff(\tpart_i,\tpart_j),\nu} = g^{ij}\frac{f_{ij}}\sqd . $$
    Recall that
    $$ g_{ij} = f_if_j + \delta_{ij} $$
    so that
    $$ g = I + (\nabla f)(\nabla f)^\top . $$
    The Sherman–Morrison formula then gives us $g^{-1}$:
    $$ g^{-1} = I - \frac{(\nabla f)(\nabla f)^\top}\den \implies g^{ij} = \delta_{ij} - \frac{f_if_j}\den . $$
    Thus
    $$ \iprod{H,\nu} = \parens{\delta^{ij} - \frac{f_if_j}\den}\frac{f_{ij}}\sqd = \frac{f_{ii}}{\sqd} - \frac{f_if_jf_{ij}}{(\den)^{3/2}} $$
    where the free indices are summed over.

    On the other hand, because $\partial_i$ forms a geodesic frame we know that
    $$ \div\parens{\frac{\nabla f}\sqd} = \partial_i\frac{f_i}{\sqd} = \frac{f_{ii}}\sqd - \frac{f_i\partial_i(\den)}{2(\den)^{3/2}} . $$
    A quick computation shows
    $$ \partial_i(\den) = 2f_jf_{ij} $$
    and so
    $$ \div\parens{\frac{\nabla f}\sqd} = \frac{f_{ii}}\sqd - \frac{f_if_jf_{ij}}{(\den)^{3/2}} = \iprod{H,\nu} $$
    as desired.

    \item In this point and the next, let
    $$ X = {\frac{\nabla f}{\sqrt{1+\norm{\nabla f}^2}}} . $$
    If $f$ is linear, then $\nabla f$ is constant and so $X$ is constant.
    Thus $\nabla X$ is zero, so its trace is zero, and thus $\div X=0$ meaning $\iprod{H,\nu}=0$.
    Since $H$ is orthogonal to $M$, this means $H=0$ so $M$ is minimal.

    \item Let us write $(x,y)$ instead of $(u_1,u_2)$ and let $r=\sqrt{x^2+y^2}$.
    Then our function is $f(r)=a\cosh^{-1}(r/a)$ whose derivative is
    $$ f'(r) = \frac1{\sqrt{r^2/a^2-1}} = \frac a{\sqrt{r^2-a^2}} $$
    Because $\nabla r=(x/r,y/r)$ we get that
    $$ \nabla f = \frac a{\sqrt{r^2-a^2}}\parens{\frac xr,\frac yr} \implies \norm{\nabla f}^2 = \frac{a^2}{r^2-a^2} . $$
    And so
    $$ \sqrt{1 + \norm{\nabla f}^2} = \frac{r}{\sqrt{r^2-a^2}} , $$
    thus
    $$ X = \frac ar\parens{\frac xr,\frac yr} = \frac a{r^2}(x,y) . $$
    Its divergence is then
    $$ \frac1a\cdot\div X = \div\parens{\frac x{r^2},\frac y{r^2}} = \pdv x\frac x{r^2} + \pdv y\frac y{r^2} . $$
    Computing shows that
    $$ \pdv x\frac x{r^2} = \pdv x\frac x{x^2+y^2} = \frac{r^2-2x^2}{r^4} $$
    so that
    $$ \frac1a\cdot\div X = \frac2{r^2} - 2\frac{x^2+y^2}{r^4} = 0  , $$
    meaning $\iprod{H,\nu}=\div X=0$ and thus $H=0$ so $M$ is minimal.
\eenum



